La segunda fase del proyecto tiene como propósito integrar equipos de red al
gestor de monitoreo Netdata instalado en la Raspberry Pi y aplicar
posteriormente los principios de gestión sobre una red MPLS. La Raspberry Pi
funcionará como el punto central de recolección de métricas, mientras que
Netdata Cloud proporcionará la interfaz centralizada para que los integrantes
del equipo accedan mediante cuentas individuales, consulten el estado de la
infraestructura y utilicen los tableros de monitoreo.

La conexión de los dispositivos dependerá de sus características. En servidores
compatibles se podrá instalar un agente de Netdata para obtener información
detallada del sistema operativo, como el uso del procesador, memoria,
almacenamiento, interfaces y servicios. En equipos de red, como switches y
routers Huawei, se utilizará el protocolo SNMP, el cual permite que un sistema
de gestión consulte variables operativas almacenadas en la base de información
de administración o MIB del dispositivo.

Para la primera parte se configurará un equipo Huawei con una dirección IP de
administración y SNMPv3. Esta versión proporciona mecanismos de autenticación
y cifrado, ofreciendo mayor seguridad que las versiones anteriores del
protocolo~\cite[sec.~1.3]{huawei-gestion-snmp}. Después se configurará Netdata para consultar el equipo desde la
Raspberry Pi y presentar métricas como el estado y tráfico de las interfaces,
errores, descartes, tiempo de actividad y utilización de recursos, según la
información que el modelo del dispositivo permita obtener mediante SNMP.

La segunda parte extenderá el procedimiento hacia una red MPLS compuesta por
routers Huawei AR651 y switches Huawei S6730H. Antes de implementar MPLS se
establecerá conectividad IP entre todos los dispositivos y se utilizará OSPF
como protocolo de enrutamiento interno. Posteriormente se configurarán y
verificarán caminos conmutados por etiquetas, inicialmente de manera estática
y después mediante LDP, observando las operaciones de inserción, intercambio y
eliminación de etiquetas realizadas por los equipos de borde y tránsito~\cite[partes~1--4]{guia-mpls-ar651-s6730h}.

De esta manera, el trabajo abarcará dos niveles de gestión. En el primero se
incorporará y supervisará un equipo Huawei individual mediante SNMPv3. En el
segundo se implementará una topología MPLS completa y se integrarán sus
dispositivos al sistema de monitoreo. Esto permitirá relacionar la
configuración lógica de la red con la observación centralizada de su
disponibilidad, rendimiento e interfaces.



